%%
%% fall-lecture10.tex
%% 
%% Made by Alex Nelson <pqnelson@gmail.com>
%% Login   <alex@lisp>
%% 
%% Started on  2026-10-20T09:22:58-0700
%% Last update 2026-10-20T09:22:58-0700
%% 

\lecture[Baire's category theorem]{}

We will need the following lemma.

\begin{lemma}[Previous Lecture]
Let $(X,\rho)$ be a complete metric space, let $\{A_{n}\}$ be a family
of closed subsets of $X$ which is monoton decreasing [so $A_{1}\supset A_{2}\supset\dots$],
and assume every $A_{n}\neq\emptyset$ is non empty and $\diam(A_{n})\to0$.
Then $\bigcap_{n\in\NN}A_{n}\neq\emptyset$.
\end{lemma}

\begin{proof}
For each $n$, pick $x_{n}\in A_{n}$ which gives us a Cauchy sequence
$(x_{n})$ which converges to some $x\in A_{n}$ for each $n\in\NN$.
\end{proof}

\begin{theorem}[Baire category]
Let $(X,\rho)$ be a complete metric space.
Let $F_{n}\subset X$ be a countable family of closed nowhere dense
subsets (so $\Interior(F_{n})=\emptyset$).
Then $\Interior(\bigcup_{n\in\NN}F_{n})=\emptyset$.
\end{theorem}

\begin{proof}
Assume for contradiction
\begin{equation}
\Interior\left(\bigcup_{n\in\NN}F_{n}\right)\neq\emptyset.
\end{equation}
Then there is a nonempty open set $O\subset\Interior(\bigcup_{n}F_{n})$.
Then there exists a point $x_{1}\in O\setminus F_{1}$ (otherwise
$O\setminus F_{1}=\emptyset$ and so $O\subset F_{1}$ which leads to
contradiction). Then there exists a small open ball
$B_{1}=B(x,r=r_{1})\Subset O\setminus F_{1}$ where $0<r_{1}<1/2$.

Then we repeat this construction, replacing $O$ with $B_{1}$ and use
$F_{2}$ now: we get $x_{2}\in B_{1}\setminus F_{2}$ and an open ball
$B_{2}=B(x_{2},r_{2})$ where $0<r_{2}<2^{-2}$ and
$B_{2}\subset B_{1}\setminus F_{2}\subset\bigcup_{n\geq3}F_{n}$.

We have a monotone decreasing sequence of open balls $\{B_{n}\}$ and
we apply our Lemma using $A_{n}=\closure{B_{n}}$. Then by our Lemma,
we get
\begin{equation}
\bigcap_{n\in\NN}\closure{B_{n}}\neq\emptyset,
\end{equation}
since $\diam(B_{n})=2\cdot2^{-n}\to0$.
Then we see this implies
\begin{subequations}
  \begin{align}
\emptyset &= O\setminus\bigcup_{n\in\NN}F_{n}\\
&=\bigcap_{n\in\NN}(O\cap F_{n})\\
&\supset\bigcap_{n\in\NN}\closure{B_{n}}\neq\emptyset,
  \end{align}
\end{subequations}
since $\closure{B_{n}}\subset O\cap F_{n}$. So we get
$\emptyset\not\subset\emptyset$ which is our contradiction.
\end{proof}

\begin{theorem}[Restatement of Baire category]
Let $(X,\rho)$ be a complete metric space. Then a countable
intersection of open dense sets is open.
\end{theorem}

\begin{remark}
The assumption is metrical, the conclusion is topological.
\end{remark}

\begin{remark}
The theorem applies to metric spaces which are not complete but are
homeomorphic to complete spaces (e.g., $S^{1}\setminus\{\mbox{pt}\}\iso\RR$
and $\RR$ is complete, so Baire's theorem applies to
$S^{1}\setminus\{\mbox{pt}\}$ despite the fact that
$S^{1}\setminus\{\mbox{pt}\}$ is not a complete metric space). So the
hypothesis really requires only ``complete metrizability'' (i.e., a
metrizable topological space whose metric is complete). But there are
\emph{even weaker versions} of Baire's theorem!
\end{remark}

\begin{definition}
We define \define{Baire Spaces} to be a topological space which
satisfies Baire's theorem.
\end{definition}

\begin{theorem}
Let $X$ be a compact Hausdorff space. Then a countable intersection of
open dense sets is dense.
\end{theorem}

(The proof is the same as before.)

\begin{remark}
Locally compact Hausdorff spaces are Baire spaces. To see this, just
use the one-point compactification.
\end{remark}

\begin{remark}[Terminology]
What's so ``categorical'' about this anyways? Well, sadly, this is
poor choice of terminology on the part of Baire.

Let $(X,\topology)$ be a topological space. Let $E\subset X$ be a
subset of $X$. We say that $E$ is \define{Category One} if
$E=\bigcup_{n\in\NN}E_{n}$ is the countable union of nowhere dense
sets $E_{n}\subset X$ [so $\Interior{\closure{E_{n}}}=\emptyset$].

We say $E$ is \define{Category Two} if it is not Category One.

This permits us to restate Baire's category theorem in the following
equivalent ways:
\begin{enumerate}
\item If $(X,\rho)$ is complete, then $X$ is category 2 in itself
\item If $(X,\rho)$ is complete, then nonempty open subsets of $X$ is category 2 in $X$
\item If $(X,\rho)$ is complete, then residual sets are dense in $X$
  (``residual'' means its complement is category 1 in $X$).
\end{enumerate}
Note that Category 1 sets are also called \define{Meager}.
\end{remark}

\begin{node}
We should think of ``residual sets'' as ``very large''.
\end{node}

There are three famous applications of Baire's category theorem.

\begin{example}[Application 1: Existence of Liouville numbers]
We say $x\in\RR$ is \define{Liouville} (and write $x\in\mathcal{L}$)
if $x\in\RR\setminus\QQ$ is irrational and for every $n\in\NN$ there
exists a $(p/q)\in\QQ$ such that
\begin{equation}
\left|x-\frac{p}{q}\right|<\frac{1}{q^{n}}.
\end{equation}
So Liouville numbers are ``irrational but almost rational''.

Compare this to algebraic numbers: $x\in\RR$ is algebraic of degree
$n\in\NN$ if there
exists a nonzero polynomial with integer coefficients $f\in\ZZ[z]$
such that $f(x)=0$. In this case, there exists some positive constant
$c>0$ such that for every $(p/q)\in\QQ$ we have
\begin{equation}
\left|x-\frac{p}{q}\right|\geq\frac{c}{q^{n}}.
\end{equation}
So this is\dots not good. And it begs the question: do Liouville
numbers even exist?
\end{example}

\begin{theorem}
The set $\mathcal{L}$ of Liouville numbers is residual in $\RR$.
But the Hausdorff dimension $\dim(\mathcal{L})=0$.
\end{theorem}

[The Hausdorff dimension is the sum of diameters of the smallest
  covering by open balls. When the Hausdorff dimension is zero, we
  should think of it as ``metrically small''. It is quite common to
  find ``large sets which are metrically small'' in analysis.]

\begin{proof}
Let $\mathcal{L}=(\RR\setminus\QQ)\cap\left(\bigcap_{n\in\NN}G_{n}\right)$
where $G_{n}=\cup_{p,q}(\frac{p}{q}-q^{-n},\frac{p}{q}+q^{-n})$. But
$G_{n}$ is open and dense, and
$\mathcal{L}^{\complement}=\QQ\cup(\bigcup_{n\in\NN}G_{n}^{\complement})$
where $G_{n}^{\complement}$ is closed nowhere-dense. Then apply
Baire's theorem. We get the result.
\end{proof}

\begin{example}[Application 2: Existence of nowhere differentiable functions]
The claim is that the  set of nowhere differentiable functions is
residual in $C[0,1]$ (continuous functions on the closed unit interval).
So almost all continuous functions are differentiable. But we can
explicitly construct the Weierstrass function
\begin{equation}
f(x)=\sum^{\infty}_{n=1}2^{-n}\cos(2^{n}x),
\end{equation}
which is not differentiable.
\end{example}

\begin{example}[Application 3: Sets of discontinuity]
Let $f\colon\RR\to\RR$. Define the set
\begin{equation}
\Dis(f)=\RR\setminus\{x\in\RR\mid\mbox{$f$ is continuous at $x$}\}.
\end{equation}
Questions:
\begin{enumerate}
\item Is there a function $f$ such that $\Dis(f)=\QQ$? [Ans: yes]
\item Can $\Dis(f)=\RR\setminus\QQ$ for some $f$? [Ans: no]
\end{enumerate}
For the first question, consider
\begin{equation}
f(x)=\begin{cases}0 & x\in\RR\setminus\QQ\\
2^{-q} & x=\frac{p}{q}\in\QQ
\end{cases}
\end{equation}
\end{example}

\begin{theorem}
If $f$ is continuous on a dense set, then $f$ is continuous on a
residual set.
\end{theorem}

\begin{example}[Application 4: what Baire was trying to do]
The class $\Baire(0)=C(\RR)$ is just continuous functions. The class
\begin{equation}
\Baire(1)=\{f\mid\exists(f_{n})\subset\Baire(0)\ldotp f_{n}\to f\mbox{ pointwise},
\forall n\in\NN\ldotp f_{n}\in\Baire(0)\}.
\end{equation}
We can inductively define the classes, for any $k\in\NN$,
\begin{equation}
\Baire(k+1)=\{f\mid\exists(f_{n})\subset\Baire(k)\ldotp f_{n}\to f\mbox{ pointwise},
\forall n\in\NN\ldotp f_{n}\in\Baire(k)\}.
\end{equation}
\end{example}

\begin{theorem}
If $f$ is $\Baire(1)$, then $\Dis(f)$ is a Category~1 set.
\end{theorem}

\begin{lemma}
If $f\colon\RR\to\RR$ is differentiable, then $f'\in\Baire(1)$ [i.e.,
  $f'$ is continuous on a residual set].
\end{lemma}

\begin{question}
Suppose we have a differentiable function $f\colon\RR\to\RR$, how can
we reconstruct $f$ from $f'(x)$ alone?

For example, if $f(x)=x^{2}\sin(1/x)$, then
\begin{equation}
f'(x)=\begin{cases}2x\sin(1/x)-\cos(1/x) & \mbox{if }x\neq0\\
0 & x=0.
\end{cases}
\end{equation}
Given this information about $f'(x)$ alone, how can we recover $f(x)$?
\end{question}